\documentclass[11pt,a4paper]{article}

\usepackage[utf8]{inputenc}
\usepackage[T1]{fontenc}
\usepackage[margin=2.2cm]{geometry}
\usepackage{parskip}
\usepackage{enumitem}
\usepackage{xcolor}
\usepackage{titlesec}
\usepackage{tcolorbox}
\usepackage{array}
\usepackage{tabularx}
\usepackage{hyperref}
\usepackage{fancyhdr}

% --- color palette (matches the website cyan/teal theme) ---
\definecolor{brand}{HTML}{0891B2}      % cyan-600
\definecolor{branddark}{HTML}{0E7490}  % cyan-700
\definecolor{brandsoft}{HTML}{ECFEFF}  % cyan-50
\definecolor{brandsoftbd}{HTML}{CFFAFE}% cyan-100
\definecolor{ink}{HTML}{1F2330}
\definecolor{muted}{HTML}{6B7280}

\hypersetup{
    colorlinks=true,
    urlcolor=branddark,
    linkcolor=branddark
}

% --- section styling ---
\titleformat{\section}
    {\Large\bfseries\color{branddark}}
    {\thesection.}{0.6em}{}
\titleformat{\subsection}
    {\normalsize\bfseries\color{ink}}
    {\thesubsection.}{0.5em}{}

\titlespacing*{\section}{0pt}{1.2em}{0.5em}
\titlespacing*{\subsection}{0pt}{0.8em}{0.3em}

\setlist[itemize]{leftmargin=1.3em, itemsep=2pt, topsep=2pt}

% --- header / footer for content pages ---
\pagestyle{fancy}
\fancyhf{}
\renewcommand{\headrulewidth}{0pt}
\fancyfoot[L]{\small\color{muted}TechnoDz -- Design Report}
\fancyfoot[R]{\small\color{muted}\thepage}

% --- handy "info" colored box ---
\newtcolorbox{infobox}[1][]{
    colback=brandsoft,
    colframe=brand,
    boxrule=0.4pt,
    arc=3pt,
    left=8pt, right=8pt, top=6pt, bottom=6pt,
    fonttitle=\bfseries\color{white},
    coltitle=white,
    colbacktitle=brand,
    title=#1
}

\begin{document}

% ==========================================================
% =====================   COVER PAGE   =====================
% ==========================================================
\thispagestyle{empty}

% top color band
\noindent\colorbox{brand}{%
    \begin{minipage}[t][3.2cm][c]{\dimexpr\textwidth-2\fboxsep\relax}
        \centering
        {\color{white}\Huge\bfseries TechnoDz}\\[0.3em]
        {\color{white}\large Mini E-Commerce Project -- Design Report}
    \end{minipage}%
}

\vspace{1.5cm}

% module + supervisor block
\begin{center}
    \begin{tcolorbox}[
        colback=brandsoft,
        colframe=brand,
        boxrule=0.5pt,
        arc=4pt,
        width=0.85\textwidth,
        left=12pt, right=12pt, top=10pt, bottom=10pt
    ]
    \centering
    {\large\bfseries\color{branddark} Module}\\[0.2em]
    {\Large PRAVAN -- Programmation Avanc\'ee}\\[1em]
    {\large\bfseries\color{branddark} Supervisor}\\[0.2em]
    {\Large Mr. BOUBENIA Mohamed}
    \end{tcolorbox}
\end{center}

\vspace{1.2cm}

% students table
\begin{center}
    {\large\bfseries\color{branddark} Project Team}
\end{center}

\vspace{0.4em}

\renewcommand{\arraystretch}{1.6}
\begin{center}
\begin{tabular}{|>{\centering\arraybackslash}m{0.5cm}|>{\centering\arraybackslash}m{7cm}|>{\centering\arraybackslash}m{4.5cm}|}
    \hline
    \rowcolor{brand}
    \textcolor{white}{\textbf{\#}} &
    \textcolor{white}{\textbf{Full Name}} &
    \textcolor{white}{\textbf{Matricule}} \\
    \hline
    1 & \underline{\hspace{6.5cm}} & \underline{\hspace{4cm}} \\
    \hline
    2 & \underline{\hspace{6.5cm}} & \underline{\hspace{4cm}} \\
    \hline
    3 & \underline{\hspace{6.5cm}} & \underline{\hspace{4cm}} \\
    \hline
\end{tabular}
\end{center}
\renewcommand{\arraystretch}{1.0}

\vspace{1.5cm}

% team number + date footer block
\begin{center}
    \begin{tabular}{ll}
        \textbf{\color{branddark}Team Number:} & \underline{\hspace{4cm}} \\[0.6em]
        \textbf{\color{branddark}Submission Date:} & \today \\
    \end{tabular}
\end{center}

\vfill

% bottom color band
\noindent\colorbox{branddark}{%
    \begin{minipage}[t][1.2cm][c]{\dimexpr\textwidth-2\fboxsep\relax}
        \centering
        {\color{white}\small Academic Year \the\year{} / \the\numexpr\year+1\relax \quad $\bullet$ \quad Mini-Project Report}
    \end{minipage}%
}

\newpage

% ==========================================================
% ===================   REPORT BODY   ======================
% ==========================================================
\setcounter{page}{1}

% short intro box at the very top
\begin{infobox}[About this report]
This document is a short design report for our mini e-commerce project.
It briefly describes the goals, the tools we used, the structure of the
code, the database, and the main features of the application.
\end{infobox}

% ----------------------------------------------------------
\section{Project Overview}

\textbf{TechnoDz} is a small online shop that sells computer parts (GPUs,
CPUs, mice, headsets, RAM, SSDs, cooling). A visitor can browse the
products, search and filter them, view a detailed product page, add items
to a cart, register an account, log in or out, and place an order. An
administrator can log in to a separate page to add, edit or delete
products and categories.

The project covers all the features asked for in the assignment: a shop
page with descriptions and prices, a search filter, a detailed product
page, a database, login and logout, a cart with cookies and sessions, and
a separate admin login.

% ----------------------------------------------------------
\section{Tools and Technologies}

We built the project on a simple LAMP-style stack on Windows with XAMPP.

\begin{itemize}
    \item \textbf{HTML} and \textbf{CSS} for the structure and the design
        (white background, cyan theme, Bootstrap Icons for the icons).
    \item \textbf{PHP} (procedural style, no framework) for the
        server-side logic.
    \item \textbf{JavaScript} for small things like the delete
        confirmation, the toast notifications, and \textbf{AJAX} on the
        cart.
    \item \textbf{MySQL} as the database, managed through phpMyAdmin.
\end{itemize}

% ----------------------------------------------------------
\section{Project Structure}

The project lives under \texttt{xampp/htdocs/nexus\_shop/}.

\begin{itemize}
    \item \textbf{Public pages} (root): \texttt{index.php},
        \texttt{product.php}, \texttt{cart.php},
        \texttt{cart\_action.php}, \texttt{checkout.php},
        \texttt{order\_confirm.php}, \texttt{login.php},
        \texttt{register.php}, \texttt{logout.php}.
    \item \textbf{Admin pages} (\texttt{/admin/}): \texttt{login.php},
        \texttt{logout.php}, \texttt{index.php} (dashboard),
        \texttt{products.php}, \texttt{product\_add.php},
        \texttt{product\_edit.php}, \texttt{product\_delete.php},
        \texttt{categories.php}.
    \item \textbf{Shared code} (\texttt{/includes/}): \texttt{config.php},
        \texttt{db.php}, \texttt{auth.php}, \texttt{cart.php},
        \texttt{helpers.php}, \texttt{header.php}, \texttt{footer.php}.
    \item \textbf{Assets}: \texttt{assets/css/style.css},
        \texttt{assets/js/app.js}, and product images in
        \texttt{assets/images/products/}.
    \item \textbf{Database scripts}: \texttt{sql/schema.sql} and
        \texttt{sql/seed.sql}.
\end{itemize}

% ----------------------------------------------------------
\section{Database Design}

The database is named \texttt{nexus\_shop} and uses InnoDB with utf8mb4.
It has seven tables:

\begin{itemize}
    \item \texttt{users} -- public accounts
        (id, username, email, password\_hash).
    \item \texttt{admins} -- staff accounts
        (id, username, password\_hash).
    \item \texttt{categories} -- product groups (id, name).
    \item \texttt{products} -- the catalog
        (id, name, description, price, image, category\_id).
    \item \texttt{cart\_items} -- saved cart for logged-in users
        (user\_id, product\_id, quantity).
    \item \texttt{orders} -- order header
        (id, user\_id, total, created\_at).
    \item \texttt{order\_items} -- order lines with a snapshot of the
        product name and unit price at checkout, so old orders stay
        correct even if the catalog changes later.
\end{itemize}

Foreign keys keep things consistent: deleting a category sets the
\texttt{category\_id} of its products to \texttt{NULL}, and deleting a
user also deletes their cart and orders.

% ----------------------------------------------------------
\section{Sessions, Cookies and Security}

\subsection{Sessions and cookies}
Both are required by the assignment, and we use them like this:
\begin{itemize}
    \item The \textbf{PHP session} stores \texttt{user\_id} for a normal
        user and \texttt{admin\_id} for an admin. The two keys are kept
        separate so an admin and a normal user cannot mix.
    \item The session also holds the active cart while the visitor is
        browsing.
    \item A \textbf{cookie} (\texttt{nexus\_cart}) saves a guest's cart
        between visits. When the guest logs in, the cookie cart is merged
        into the user's saved cart in the database, and the cookie is
        cleared.
\end{itemize}

\subsection{Security}
\begin{itemize}
    \item Every SQL query uses \textbf{prepared statements}
        (mysqli with \texttt{bind\_param}). User input is never put
        directly into SQL.
    \item Passwords are stored with \texttt{password\_hash()} (bcrypt)
        and checked with \texttt{password\_verify()}.
    \item All output goes through an \texttt{h()} helper that calls
        \texttt{htmlspecialchars()} to prevent XSS.
    \item Admin pages call \texttt{require\_admin()} at the top to
        redirect anyone who is not logged in as an admin.
\end{itemize}

% ----------------------------------------------------------
\section{Main Features}

\subsection{Shop and search}
The home page (\texttt{index.php}) shows all products in a grid. The
header has a search bar with a magnifier icon. Below the hero banner
there is a strip of category chips (\textit{All}, GPU, CPU, Mouse, ...).
Clicking a chip filters the grid; the active chip is highlighted. The
search and the category filter work together.

\subsection{Product detail}
Clicking a card opens \texttt{product.php?id=N} which shows the full
description, the price and an \emph{Add to cart} button. If the id does
not exist, the page shows a clear ``Product not found'' message.

\subsection{Cart}
The cart has three layers:
\begin{itemize}
    \item Session: the working cart during the visit.
    \item Cookie: keeps a guest cart between visits.
    \item Database (\texttt{cart\_items}): saves the cart for logged-in
        users across devices.
\end{itemize}
Prices are shown in \textbf{DZD} (we store the value once and multiply by
260 in the \texttt{money()} helper).

\subsection{Checkout and order}
Only logged-in users can place an order. \texttt{checkout.php} re-reads
the products from the database, computes the total, opens a transaction,
inserts one row in \texttt{orders} and one row in \texttt{order\_items}
per cart line, commits, then empties the cart and shows
\texttt{order\_confirm.php} with the order id.

\subsection{Authentication}
Users register with a username, email and password (validated and
hashed). Login uses the email; logout clears the session and the guest
cart cookie. The admin uses a separate login at \texttt{/admin/login.php}
with a username and password (default seeded admin: \texttt{admin /
admin123}).

\subsection{Admin panel}
The admin dashboard shows simple counts (products, categories, orders)
and quick links. The admin can:
\begin{itemize}
    \item Add a product (with an image upload checked by MIME type, and a
        positive price check).
    \item Edit a product (with an optional new image).
    \item Delete a product (POST-only, with a JavaScript confirm dialog).
    \item Add or delete categories.
\end{itemize}

\subsection{AJAX}
To make the site feel smooth, \emph{Add to cart}, \emph{Update quantity}
and \emph{Remove} on the cart page use AJAX (\texttt{fetch} in vanilla
JavaScript). The handler \texttt{cart\_action.php} returns JSON when the
request comes from JavaScript, so the cart badge in the header and the
line totals update without reloading the page. If JavaScript is off, the
same forms still work the classic way.

% ----------------------------------------------------------
\section{Conclusion}

The project shows the main steps of building a small e-commerce site:
a clean database, separated public and admin areas, sessions and cookies
for the cart, prepared statements for security, and a simple but pleasant
interface. The code is kept simple on purpose so it stays easy to read
and to demo.

\vspace{0.5em}
\begin{infobox}[Thank you]
We thank Mr. BOUBENIA Mohamed for the guidance and the feedback during
the PRAVAN module.
\end{infobox}

\end{document}
